% --- tabla_contenido.tex ---

% 1. Configuración local: FUENTES, SIN PUNTOS y NUMERACIÓN OPCIONAL
% El segundo bloque de llaves {} maneja la numeración. 

% PARTES: Si quieres número, cambia el {} por {Parte \thecontentslabel\quad}
\titlecontents{part}[-1.1em]
{\bfseries\sbseries\normalsize}
{\thecontentslabel\quad} % <--- Aquí aparece el número si existe
{}
{\hfill\small\lhseries\contentspage\vspace{1.0mm}}[]

% CAPÍTULOS: Cambia {\thecontentslabel.\quad} por {} si quieres ocultar el número
\titlecontents{chapter}[0em]
{\medseries\normalsize}
{\thecontentslabel.\quad} % <--- Muestra "1. " o "2. "
{}
{\hfill\small\lhseries\contentspage\vspace{2.5mm}}[]

% SECCIONES: Cambia {\thecontentslabel\quad} por {} si quieres ocultar el número
\titlecontents{section}[1.5em]
{\lhseries}
{\thecontentslabel\quad} % <--- Muestra "1.1 "
{}
{\hfill\small\lhseries\contentspage\vspace{0.5mm}}[]

% SUBSECCIONES
\titlecontents{subsection}[3em]
{\lhseries}
{\thecontentslabel\quad}
{}
{\hfill\small\lhseries\contentspage\vspace{0.5mm}}[]

% 2. Ocultar numeración de página física del índice
\cleardoublepage
\pagestyle{empty} 
\assignpagestyle{\chapter}{empty} 

% 3. Título del índice
\renewcommand*\contentsname{\lhseries Contenido}

% 4. Generación del índice
\tableofcontents

% 5. Restauración para el resto del libro
\cleardoublepage
\pagestyle{esstyle}